\documentclass{article}
\usepackage{amsmath,amssymb,amsthm}
\usepackage{algorithm}
\usepackage{algpseudocode}
\usepackage{geometry}
\geometry{margin=1in}
\newtheorem{theorem}{Theorem}
\newtheorem{lemma}{Lemma}
\newtheorem{assumption}{Assumption}
\newtheorem{corollary}{Corollary}
\newcommand{\R}{\mathbb{R}}
\newcommand{\Cset}{\mathcal{C}}
\newcommand{\grad}{\nabla}
\newcommand{\gap}{\mathcal{G}}
\newcommand{\opt}{\star}
\newcommand{\inner}[2]{\langle #1,#2\rangle}
\begin{document}

\section*{Algorithm Name}
\textbf{Conditional Gradient (Frank–Wolfe) for Non‑Convex Smooth Optimization}

\section*{Mathematical Formulation}
Let $f:\R^{d}\rightarrow\R$ be differentiable and $L$–smooth on a compact convex set 
$\Cset\subset\R^{d}$ (i.e., $\| \grad f(x)-\grad f(y)\|\le L\|x-y\|$ for all $x,y\in\Cset$).
Given an initial point $x_{0}\in\Cset$ and a stepsize $\gamma\in(0,1]$, the algorithm generates
\[
\begin{aligned}
s_{k} &:= \arg\min_{s\in\Cset}\inner{\grad f(x_{k})}{s}, \tag{1}\\
x_{k+1} &:= (1-\gamma)x_{k}+\gamma s_{k}. \tag{2}
\end{aligned}
\]
The \emph{Frank–Wolfe (FW) gap} at $x$ is
\[
\gap(x):=\max_{s\in\Cset}\inner{\grad f(x)}{x-s}
          =\inner{\grad f(x)}{x-s^{\mathrm{FW}}(x)},
\]
where $s^{\mathrm{FW}}(x)$ denotes the solution of (1).  Note that $\gap(x)\ge0$ and
$\gap(x)=0$ iff $x$ is a stationary point of the constrained problem
$\min_{x\in\Cset}f(x)$.

\section*{Pseudocode}
\begin{algorithm}[H]
\caption{Conditional Gradient (Frank–Wolfe) for Non‑Convex $f$}
\begin{algorithmic}[1]
\Require Compact convex set $\Cset$, smooth $f$, stepsize $\gamma\in(0,1]$, max iter $T$
\State Choose $x_{0}\in\Cset$
\For{$k=0,1,\dots,T-1$}
    \State Compute gradient $g_{k}\gets\grad f(x_{k})$
    \State \textbf{Linear minimization oracle:}
          \[
          s_{k}\gets\arg\min_{s\in\Cset}\inner{g_{k}}{s}
          \tag{Line~\ref{linoracle}}
          \]
    \State Update 
          \[
          x_{k+1}\gets (1-\gamma)x_{k}+\gamma s_{k}
          \tag{Line~\ref{update}}
          \]
\EndFor
\State \Return $\{x_{k}\}_{k=0}^{T}$
\end{algorithmic}
\end{algorithm}

\section*{Dry Run Example}
Consider the one‑dimensional problem
\[
f(w)=(w-3)^{2},\qquad \Cset=[0,5],\qquad \gamma=0.1,\qquad w_{0}=0.
\]
The gradient is $\grad f(w)=2(w-3)$.  The linear minimization oracle solves
\[
s_{k}=\arg\min_{s\in[0,5]} 2(w_{k}-3)s=
\begin{cases}
5,& w_{k}<3,\\
0,& w_{k}>3,
\end{cases}
\]
(because the coefficient $2(w_{k}-3)$ is negative when $w_{k}<3$).

\begin{center}
\begin{tabular}{c|c|c|c|c}
$k$ & $w_{k}$ & $\grad f(w_{k})$ & $s_{k}$ & $w_{k+1}= (1-\gamma)w_{k}+\gamma s_{k}$ \\ \hline
0 & $0$ & $-6$ & $5$ & $0.9\cdot0+0.1\cdot5=0.5$ \\
1 & $0.5$ & $-5$ & $5$ & $0.9\cdot0.5+0.1\cdot5=0.95$ \\
2 & $0.95$ & $-4.1$ & $5$ & $0.9\cdot0.95+0.1\cdot5=1.355$\\
\end{tabular}
\end{center}
The corresponding objective values are $f(0)=9$, $f(0.5)=6.25$, $f(0.95)=4.2025$, $f(1.355)=2.949\ldots$, showing monotone decrease.

\newpage
\section*{Assumptions}
\begin{assumption}[Smoothness]\label{ass:smooth}
$f$ is $L$‑smooth on $\Cset$: for all $x,y\in\Cset$,
\[
f(y)\le f(x)+\inner{\grad f(x)}{y-x}+\frac{L}{2}\|y-x\|^{2}.
\tag{A1}
\]
\end{assumption}
\begin{assumption}[Compactness]\label{ass:compact}
$\Cset$ is convex and has finite diameter 
\[
D:=\max_{x,y\in\Cset}\|x-y\|<\infty .
\tag{A2}
\]
\end{assumption}
\begin{assumption}[Bounded suboptimality]\label{ass:opt}
Denote $f^{*}:=\min_{x\in\Cset}f(x)$.  The initial suboptimality is finite:
\[
\Delta_{0}:=f(x_{0})-f^{*}<\infty .
\tag{A3}
\]
\end{assumption}

\section*{Main Convergence Theorem}
\begin{theorem}[Convergence of the non‑convex Frank–Wolfe method]\label{thm:fw}
Under Assumptions \ref{ass:smooth}–\ref{ass:opt} and a constant stepsize 
$\gamma\in(0,1]$, the iterates $\{x_{k}\}$ generated by \eqref{linoracle}–\eqref{update}
satisfy for any integer $T\ge1$:
\[
\min_{0\le k<T}\gap(x_{k})
\;\le\;
\frac{f(x_{0})-f^{*}}{\gamma T}
\;+\;
\frac{\gamma L D^{2}}{2}.
\tag{3}
\]
Consequently, choosing $\gamma=\sqrt{\frac{2\Delta_{0}}{L D^{2}T}}$ yields 
\[
\min_{0\le k<T}\gap(x_{k})\le
\mathcal{O}\!\left(\frac{1}{\sqrt{T}}\right).
\]
\end{theorem}

\section*{Theoretical Properties}
\begin{itemize}
\item \textbf{Stability.}  The bound \eqref{3} depends only on $L$, $D$, the initial gap
$\Delta_{0}$ and the stepsize $\gamma$.  Hence the method is robust to the choice of
initial point as long as $x_{0}\in\Cset$.
\item \textbf{Hyperparameter Sensitivity.}  A larger $\gamma$ accelerates the linear term
$\frac{f(x_{0})-f^{*}}{\gamma T}$ but inflates the quadratic term $\frac{\gamma L D^{2}}{2}$.
The optimal constant stepsize that balances the two terms is 
$\gamma^{\star}= \sqrt{\frac{2\Delta_{0}}{L D^{2}T}}$, yielding the $O(1/\sqrt{T})$ rate.
\item \textbf{Comparison to Gradient Descent.}  For constrained non‑convex problems,
projected gradient descent typically achieves the same $O(1/\sqrt{T})$ rate
but requires a projection onto $\Cset$ at each iteration.
The Frank–Wolfe method replaces the projection by a linear minimization oracle,
which can be cheaper when $\Cset$ has a simple structure (e.g., simplex, trace‑norm ball).
\end{itemize}

\section*{Discussion}
The theorem guarantees that at least one iterate among the first $T$ steps
has a small FW gap, which is a stationarity measure for constrained problems.
Because the algorithm never leaves $\Cset$, feasibility is maintained automatically.
If $f$ happens to be convex, the same bound implies the standard $O(1/T)$ primal
convergence after averaging the iterates.  In the non‑convex case, the $O(1/\sqrt{T})$
rate matches the best known guarantees for first‑order methods that only use
oracle information.

\newpage
\section*{Preliminaries (Definitions and Lemmas)}
\begin{definition}[FW gap]
For $x\in\Cset$,
\[
\gap(x):=\max_{s\in\Cset}\inner{\grad f(x)}{x-s}
       =\inner{\grad f(x)}{x-s^{\mathrm{FW}}(x)},
\]
where $s^{\mathrm{FW}}(x)$ solves the linear oracle \eqref{linoracle}.
\end{definition}
\begin{lemma}[Descent Lemma for $L$‑smooth functions]\label{lem:descent}
For any $x,y\in\Cset$,
\[
f(y)\le f(x)+\inner{\grad f(x)}{y-x}
          +\frac{L}{2}\|y-x\|^{2}.
\tag{L1}
\]
\end{lemma}
\begin{proof}
Standard consequence of $L$‑smoothness; see Assumption \ref{ass:smooth}.
\end{proof}
\begin{lemma}[Bound on the step]\label{lem:stepbound}
For the update \eqref{update},
\[
\|x_{k+1}-x_{k}\|\le \gamma D .
\tag{L2}
\]
\end{lemma}
\begin{proof}
Since $s_{k},x_{k}\in\Cset$, $\|s_{k}-x_{k}\|\le D$.  Using \eqref{update},
\[
x_{k+1}-x_{k}= \gamma (s_{k}-x_{k})\;\Rightarrow\;
\|x_{k+1}-x_{k}\|=\gamma\|s_{k}-x_{k}\|\le\gamma D .
\]
\end{proof}

\newpage
\section*{Main Proof (Step‑by‑Step Derivation)}
We prove Theorem \ref{thm:fw}.  Throughout the proof we fix an arbitrary $k\ge0$.

\begin{enumerate}
\item[Step 1.] Apply Lemma \ref{lem:descent} with $x=x_{k}$ and $y=x_{k+1}$:
\[
f(x_{k+1})\le f(x_{k})+\inner{\grad f(x_{k})}{x_{k+1}-x_{k}}
               +\frac{L}{2}\|x_{k+1}-x_{k}\|^{2}.
\tag{4}
\]

\item[Step 2.] Substitute the update rule \eqref{update} into the inner product:
\[
x_{k+1}-x_{k}= \gamma(s_{k}-x_{k})\;\Rightarrow\;
\inner{\grad f(x_{k})}{x_{k+1}-x_{k}}
   =\gamma\inner{\grad f(x_{k})}{s_{k}-x_{k}}.
\tag{5}
\]

\item[Step 3.] By definition of the FW gap,
\[
\inner{\grad f(x_{k})}{s_{k}-x_{k}}
   = -\gap(x_{k}).
\tag{6}
\]

\item[Step 4.] Use Lemma \ref{lem:stepbound} to bound the quadratic term:
\[
\|x_{k+1}-x_{k}\|^{2}= \gamma^{2}\|s_{k}-x_{k}\|^{2}
                    \le \gamma^{2} D^{2}.
\tag{7}
\]

\item[Step 5.] Combine (4)–(7):
\[
\begin{aligned}
f(x_{k+1})
&\le f(x_{k}) + \gamma\bigl(-\gap(x_{k})\bigr)
   +\frac{L}{2}\gamma^{2}D^{2}\\[2mm]
&= f(x_{k}) -\gamma\gap(x_{k}) +\frac{\gamma^{2}L D^{2}}{2}.
\end{aligned}
\tag{8}
\]

\item[Step 6.] Rearrange (8) to isolate the gap term:
\[
\gamma\gap(x_{k})
\le f(x_{k})-f(x_{k+1}) +\frac{\gamma^{2}L D^{2}}{2}.
\tag{9}
\]

\item[Step 7.] Sum inequality (9) over $k=0,\dots,T-1$:
\[
\gamma\sum_{k=0}^{T-1}\gap(x_{k})
\le \sum_{k=0}^{T-1}\bigl[f(x_{k})-f(x_{k+1})\bigr]
   +\frac{T\gamma^{2}L D^{2}}{2}.
\tag{10}
\]

\item[Step 8.] The telescoping sum simplifies:
\[
\sum_{k=0}^{T-1}\bigl[f(x_{k})-f(x_{k+1})\bigr]
= f(x_{0})-f(x_{T})\le f(x_{0})-f^{*}.
\tag{11}
\]

\item[Step 9.] Insert (11) into (10) and divide by $\gamma T$:
\[
\frac{1}{T}\sum_{k=0}^{T-1}\gap(x_{k})
\le \frac{f(x_{0})-f^{*}}{\gamma T}
   +\frac{\gamma L D^{2}}{2}.
\tag{12}
\]

\item[Step 10.] Since the minimum of a set does not exceed its average,
\[
\min_{0\le k<T}\gap(x_{k})
\le\frac{1}{T}\sum_{k=0}^{T-1}\gap(x_{k}).
\tag{13}
\]

\item[Step 11.] Combine (12) and (13) to obtain the claimed bound:
\[
\boxed{\;
\min_{0\le k<T}\gap(x_{k})
\le \frac{f(x_{0})-f^{*}}{\gamma T}
   +\frac{\gamma L D^{2}}{2}\;}
\tag{14}
\]
which is exactly \eqref{3}.  $\square$
\end{enumerate}

\section*{Rate Derivation (With Constants)}
From \eqref{14}, set $\gamma=\sqrt{\tfrac{2\Delta_{0}}{L D^{2}T}}$ where $\Delta_{0}=f(x_{0})-f^{*}$.
Then
\[
\frac{\Delta_{0}}{\gamma T}
   =\frac{\Delta_{0}}{T}\sqrt{\frac{L D^{2}T}{2\Delta_{0}}}
   =\sqrt{\frac{L D^{2}\Delta_{0}}{2T}},
\qquad
\frac{\gamma L D^{2}}{2}
   =\frac{L D^{2}}{2}\sqrt{\frac{2\Delta_{0}}{L D^{2}T}}
   =\sqrt{\frac{L D^{2}\Delta_{0}}{2T}}.
\]
Hence
\[
\min_{0\le k<T}\gap(x_{k})
\le 2\sqrt{\frac{L D^{2}\Delta_{0}}{2T}}
= \mathcal{O}\!\left(\frac{1}{\sqrt{T}}\right).
\]
Thus the algorithm attains the optimal $O(1/\sqrt{T})$ stationarity rate for
smooth non‑convex constrained optimization.

\section*{Special Cases}
\begin{itemize}
\item \textbf{Convex $f$.}  When $f$ is convex, the FW gap upper‑bounds the primal
suboptimality:
\[
f(x_{k})-f^{*}\le\gap(x_{k}),
\]
so \eqref{14} directly yields the classical $O(1/T)$ convergence of the
Frank–Wolfe method after averaging.
\item \textbf{Exact line search.}  If instead of a constant $\gamma$ we choose
$\gamma_{k}=\arg\min_{\gamma\in[0,1]} f\bigl((1-\gamma)x_{k}+\gamma s_{k}\bigr)$,
the quadratic term in \eqref{8} disappears, leading to the sharper bound
$\min_{k<T}\gap(x_{k})\le \frac{2\Delta_{0}}{T}$.
\item \textbf{Polytope $\Cset$ with bounded curvature.}  When $\Cset$ is a polytope,
the diameter $D$ can be replaced by the curvature constant $C_{f}$,
yielding a bound of the same form with $L D^{2}$ replaced by $C_{f}$.
\end{itemize}

\end{document}